\documentclass[11pt]{article}

\usepackage[a4paper]{geometry}
\geometry{left=2.0cm,right=2.0cm,top=2.5cm,bottom=2.5cm}

\usepackage{ctex}
\usepackage{amsmath,amsfonts,graphicx,amssymb,bm,amsthm}
\usepackage{algorithm,algorithmicx}
\usepackage[noend]{algpseudocode}
\usepackage{fancyhdr}
\usepackage{amssymb}
\usepackage{listings}
\usepackage{xcolor}
\usepackage{enumitem}
\usepackage{url}

\newlist{problem}{enumerate}{1}
\setlist[problem]{
    label={},
    leftmargin=4em,      % 自动计算左边距以适应编号宽度
    align=left,        % 编号左对齐
    labelsep=2em,    % 编号与正文的间距
}

\newlist{solution}{enumerate}{1}
\setlist[solution]{
    label={\textbf{解}},
    leftmargin=4em,      % 自动计算左边距以适应编号宽度
    align=left,        % 编号左对齐
    labelsep=2em,    % 编号与正文的间距
}

\newlist{prove}{enumerate}{1}
\setlist[prove]{
    label={\textbf{证明}},
    leftmargin=4em,      % 自动计算左边距以适应编号宽度
    align=left,        % 编号左对齐
    labelsep=2em,    % 编号与正文的间距
}

% 段落间距
\setlength{\parskip}{1em}
\allowdisplaybreaks

\pagestyle{fancy}
\lhead{\kaishu 中国科学院大学}
\chead{}
\rhead{\kaishu 2026年春季学期汇编语言}

    
\begin{document}

% 标题区域
\begin{center}
    {\LARGE \bf 汇编语言大作业2 \quad 实验报告} \\
\end{center}
\begin{kaishu}
    小组成员：吴竞宸、黄奕皓、张晨轩 \hfill \quad 2026/05/26
\end{kaishu}

\section{概述}
本次实验采用小组合作的方式完成，开发环境为 Linux，成员通过微信群交流开发进度，使用 Github 共同管理代码。

Github 仓库网址为：\url{https://github.com/Dandelom/UCAS-ASM-PRJ2}

实验要求实现4096*4096单精度浮点的矩阵转置程序，在实现简单转置的基础上，我们采用了4种优化方法，分别为：
\begin{enumerate}
    \item 分块转置（64*64）
    \item 4x4 SIMD 转置
    \item 8x8 AVX 转置
    \item AVX 转置 + Prefetch
\end{enumerate}

其中，分块转置以及最先实现的简单转置使用 x86-32 指令集，而后三种优化方式使用 x86-64 指令集完成，
分别由 main32.c 和 main64.c 控制。

建议的编译指令：

gcc -m32 main32.c simply\_transpose.s blockwise\_transpose.c diagonal\_transpose.s off\_diagonal\_\allowbreak transpose.s -o ../bin/transpose32

gcc main64.c simd\_block\_transpose.s simd\_transpose.c avx\_block\_transpose.s avx\_transpose.c avx\_prefetch.c avx\_dual\_prefetch\_transpose.s avx\_prefetch\_transpose.s -o ../bin/transpose64

src 文件夹下还有 generate.c 用于生成待转置的 4096*4096 单精度浮点矩阵，采用二进制文件形式存储；check.c 用于检查转置是否正确。测试数据和输出数据（.bin 文件）存放在 tests 文件夹下。

建议的编译指令：

gcc generate.c -o ../bin/generate

gcc check.c -o ../bin/check

\section{实验结果}
\subsection{实验数据}
程序运行环境为 Windows 11 下 VMware Workstation 内的 Ubuntu 24.04.3 虚拟机。为虚拟机
分配的内存为 6 GB，处理器内核总数为 4。物理机的 CPU 型号为 Intel(R) Core(TM) i9-14900HX (2.20 GHz)。

为简单转置及各加速方法采用5次运行取平均值的方法，原始数据记录如下。
\begin{table}[H]
    \centering
    单位：ms \\
    \begin{tabular}{|c|c|c|c|c|c|}\hline
         & 简单转置 & 分块转置 & SIMD & AVX & AVX + Prefetch \\ \hline
        1 & 72.81 & 38.35 & 53.50 & 22.69 & 11.69 \\ \hline
        2 & 70.73 & 37.25 & 50.87 & 24.38 & 11.80 \\ \hline
        3 & 76.10 & 38.19 & 48.18 & 24.29 & 11.46 \\ \hline
        4 & 66.00 & 38.85 & 48.86 & 23.52 & 14.11 \\ \hline
        5 & 69.96 & 39.56 & 49.35 & 27.69 & 13.41 \\ \hline
    \end{tabular}
    \caption{原始数据}
\end{table}

计算得到最终算法时间与加速比如下：
\begin{table}[H]
    \centering
    \begin{tabular}{ccc}\hline
        算法 & 运行时间(ms) & 加速比 \\ \hline
        简单转置 & 71.12 & 1.00 \\
        分块转置 & 38.44 & 1.85 \\
        SIMD & 50.15 & 1.42 \\
        AVX & 24.51 & 2.90 \\
        AVX + Prefetch & 12.49 & 5.69 \\ \hline
    \end{tabular}
    \caption{算法时间与加速比}
\end{table}

\subsection{算法分析}
普通转置在处理行和列时，由于大矩阵（$4096 \times 4096$）的单行数据量
远超 L1/L2 缓存容量，按列读取或写入时会产生巨大的跨步访问（Stride Visit），
导致严重的 Cache Line 缺失（Cache Miss）。分块转置将大矩阵切分为局部的小块，
使得小块的数据可以完全容纳在高速缓存（L1 Cache）中，
极大地提高了时间局部性和空间局部性，加速比达到了1.85。

从数据上看，4x4 SIMD 的加速比低于普通分块转置。分析代码，可能的原因是频繁的 memcpy 带来的数据吞吐损耗抵消了 SIMD 节省下来的计算时间。

而 8x8 AVX 转置在分块的基础上利用了 SIMD（单指令多数据） 寄存器。AVX 拥有 256 位的 \%ymm 寄存器，单条指令可以同时处理 8 个单精度浮点数。它的核心优势是用硬件寄存器内的洗牌指令（Shuffle/Unpack/Permute）替代了传统的标量交换循环，消除了分支预测失败的开销，并最大化了指令级并行度（ILP）。加速比达到2.90。

AVX + Prefetch 转置是性能发生质飞跃的一步。即使数据在 L1 缓存里处理得再快，当跳到下一个物理上不连续的内存块时，CPU 仍然需要停下来等待总线从主存中加载数据。通过在 C 语言端计算下一轮循环的地址并传入汇编，利用 prefetcht0 指令在当前块执行洗牌计算的同时，异步将下一个块的数据提前拉入 L1 缓存。这就完美掩盖了内存访问延迟（Memory Latency），使 CPU 核心几乎不用处于空转（Stall）状态。
这使得加速比来到了最高的5.69。

\end{document}